\documentclass[10pt]{article}
\usepackage[margin=0.9in]{geometry}
\usepackage[T1]{fontenc}
\usepackage{lmodern}
\usepackage{booktabs}
\usepackage{tabularx}
\usepackage{xltabular}
\usepackage{array}
\usepackage{pdflscape}
\usepackage{xcolor}
\usepackage{listings}
\usepackage[hyphens]{url}
\usepackage[colorlinks=true,linkcolor=blue!50!black,urlcolor=blue!50!black]{hyperref}
\usepackage{enumitem}

% file:line references; \path allows line breaks inside long paths
\newcommand{\f}[1]{\path{#1}}
\newcommand{\yes}{\textbf{yes}}
\newcommand{\no}{\textbf{no}}
\newcommand{\partl}{\textbf{partial}}

\lstset{basicstyle=\ttfamily\scriptsize,breaklines=true,frame=single,
  columns=fullflexible,keepspaces=true,showstringspaces=false,
  backgroundcolor=\color{black!3}}

\setlist{itemsep=2pt,topsep=3pt}

\title{Attack Surface for a Quantization-Conditioned RAG Index Backdoor:\\
A Code-Level Audit of FAISS, LangChain, LlamaIndex, Milvus, Qdrant, and Weaviate}
\author{Fall 2026 Conversion-Integrity Project}
\date{27 September 2026}

\begin{document}
\maketitle

\noindent\textbf{Snapshot audited.} All paths are relative to \f{RAG-Related/}.
Commits: faiss \texttt{fdb9535c1}, langchain \texttt{d7508191bc},
llama\_index \texttt{169e450aa}, milvus \texttt{acd8a6770e}, qdrant \texttt{6ab21cac1},
weaviate \texttt{e3f4df7ba8}. Line numbers refer to these commits. ``Not found'' means we
searched the repository and found no such code. It does not cover code outside the repository.

%==========================================================================
\section{Verdict}
%==========================================================================

The four-part surface (Q1--Q4) exists, but not in the form ``RAG stacks ship PQ by
default''. That claim is \textbf{false} for five of the six repositories. \textbf{Q1:} PQ/SQ
codebook construction exists in FAISS (C++ k-means per sub-vector), Qdrant (own Rust PQ/SQ),
and Weaviate (own Go PQ/SQ/RQ). Milvus exposes PQ/SQ index types but delegates the
implementation to the external \texttt{knowhere} library, which is not in the repo. LangChain
and LlamaIndex implement no quantization and delegate to a backend. \textbf{Q2:} PQ is opt-in
everywhere. Scalar quantization is on by default in exactly one place: \textbf{Milvus}. There,
any vector index created without an explicit \texttt{index\_type}, including the SDK's
quick-setup path, resolves to \texttt{HNSW\_SQ} with \texttt{sq\_type=SQ4U} and an FP16
refine stage. Weaviate and Qdrant leave compression off unless the collection or server is
configured for it. Once PQ is enabled in Weaviate, it trains automatically when a shard
passes \texttt{trainingLimit}=100{,}000 vectors (a scale trigger). \textbf{Q3:} every system
has a shippable compressed artifact. The artifacts are: a FAISS \texttt{write\_index} file
(codebooks and codes, with no raw vectors for IVFPQ); a LlamaIndex persist directory (which
removes embeddings from the docstore, so the FAISS file is the only copy of the vectors); a
Qdrant snapshot tarball (\texttt{quantized.meta.json} holds the centroids as plain JSON);
Weaviate backups (HNSW commit logs carry the PQ encoders); and Milvus snapshots, including
\texttt{RestoreExternalSnapshot}, which copies prebuilt index files. \textbf{Q4:} no
repository checks a loaded quantized index against the source embeddings. The only checks we
found are structural (shape and size), identity-based (Milvus build-ID ownership), or
optional whole-file SHA-256 values supplied by the same party that ships the file (Qdrant).
The real limit on the attack is \emph{rescoring}. Weaviate re-ranks compressed candidates
with full-precision vectors by default, and Milvus's default index has an FP16 refine stage.
Both reduce the attack to candidate-set suppression. Qdrant's PQ/SQ do \emph{not} rescore by
default, and FAISS IVFPQ never does unless the index is wrapped in \texttt{IndexRefine}.
\textbf{Strongest:} FAISS IVFPQ loaded through LlamaIndex's
\texttt{FaissVectorStore.from\_persist\_dir} (Section~\ref{sec:strongest}). Next come Qdrant
with PQ enabled (no default rescoring, snapshot recovery from a URL) and Milvus's default
\texttt{HNSW\_SQ} (quantized by default, but with refine). \textbf{Weakest:} LangChain core,
whose in-repo stores are exact and whose FAISS wrapper is not in this repo, and Weaviate,
which rescores by default. The honest threat model is therefore an \emph{explicitly
configured or scale-triggered} deployment, with Milvus SQ as the one default-on case, rather
than a default-on PQ supply chain.

%==========================================================================
\begin{landscape}
\section{Summary table}
%==========================================================================
{\scriptsize
\renewcommand{\arraystretch}{1.25}
\begin{xltabular}{\linewidth}{@{}>{\raggedright\arraybackslash}p{1.35cm}
  *{6}{>{\raggedright\arraybackslash}X}@{}}
\toprule
\textbf{Repo} & \textbf{Q1 PQ/SQ exists} & \textbf{Q2 default vs.\ opt-in (exact default)} &
\textbf{Q3 shippable compressed artifact} & \textbf{Q4 verification on load} &
\textbf{Q5 seed / determinism} & \textbf{Q6 direct codebook override} \\
\midrule\endhead
\midrule
FAISS &
\yes. \f{faiss/faiss/impl/ProductQuantizer.cpp:130-206} (k-means per sub-vector);
\f{faiss/faiss/IndexIVFPQ.cpp:76-95}; OPQ \f{faiss/faiss/VectorTransform.cpp:1246}. &
\textbf{opt-in}. There is no default index: \texttt{index\_factory} needs a description
(\f{faiss/faiss/index_factory.h:17-20}); \texttt{"Flat"} gives exact \texttt{IndexFlat}
(\f{faiss/faiss/index_factory.cpp:809-811}). &
\yes. \texttt{write\_index} stores centroids and codes
(\f{faiss/faiss/impl/index_write.cpp:185-190,801-810}). &
\no. Shape checks only (\f{faiss/faiss/impl/index_read.cpp:735-768}). &
\textbf{seeded}: \texttt{seed=1234} (\f{faiss/faiss/Clustering.h:47-50}); OPQ 1234
(\f{faiss/faiss/VectorTransform.cpp:1262}). Seed and \texttt{cp} are \emph{not serialized}. &
\yes. \texttt{set\_params} (\f{faiss/faiss/impl/ProductQuantizer.cpp:76-80}),
\texttt{Train\_hot\_start}, \texttt{add\_sa\_codes} (\f{faiss/faiss/Index.h:430}),
\texttt{update\_entries}. \\
\midrule
LangChain &
\no\ in repo. FAISS is a lazy import from \texttt{langchain\_community}, which is not
vendored (\f{langchain/libs/langchain/langchain_classic/vectorstores/faiss.py:11}). &
\textbf{no quantization}. \texttt{InMemoryVectorStore} is exact cosine
(\f{langchain/libs/core/langchain_core/vectorstores/in_memory.py:314}); the Qdrant partner
creates plain \texttt{VectorParams} (\f{langchain/libs/partners/qdrant/langchain_qdrant/qdrant.py:944-982}). &
\partl. \texttt{InMemoryVector\-Store}: \texttt{dump}/\texttt{load} is full-precision JSON
(\f{langchain/libs/core/langchain_core/vectorstores/in_memory.py:517-546});
\texttt{from\_existing\_collection} attaches to a pre-built Qdrant collection. &
\no. The Qdrant partner validates only size and distance
(\f{langchain/libs/partners/qdrant/langchain_qdrant/qdrant.py:1230,1240}). &
n/a (no training code). &
n/a in repo. \\
\midrule
LlamaIndex &
\no\ own PQ. It wraps any \texttt{faiss.Index}
(\f{llama_index/llama-index-integrations/vector_stores/llama-index-vector-stores-faiss/llama_index/vector_stores/faiss/base.py:65-82}). &
\textbf{no quantization by default}: \texttt{SimpleVectorStore} (exact)
(\f{llama_index/llama-index-core/llama_index/core/storage/storage_context.py:105-107}).
The Milvus integration forces \texttt{FLAT} (\f{.../milvus/base.py:1406}). &
\yes. \texttt{persist} calls \texttt{faiss.write\_index} (\f{.../faiss/base.py:171}), and
the docstore \emph{removes} embeddings (\f{llama_index/llama-index-core/llama_index/core/indices/vector_store/base.py:196-203}). &
\no. \texttt{from\_persist\_path} calls \texttt{faiss.read\_index} and trusts the result
(\f{.../faiss/base.py:100-117}). &
Inherits FAISS. &
\yes, inherited: any user-supplied FAISS object is accepted. \\
\midrule
Milvus &
\partl. Types \texttt{IVF\_PQ}, \texttt{HNSW\_PQ}, \texttt{HNSW\_SQ}, \texttt{SCANN}
(\f{milvus/client/index/common.go:46-57}); the implementation is in external \texttt{knowhere}
(\f{milvus/internal/core/thirdparty/knowhere/CMakeLists.txt:17-18}). &
\textbf{SQ on by default}. AUTOINDEX resolves to
\texttt{HNSW\_SQ}/\texttt{SQ4U}/refine \texttt{FP16}
(\f{milvus/pkg/util/paramtable/autoindex_param_nocuda.go:26}) when no
\texttt{index\_type} is given (\f{milvus/internal/proxy/task_index.go:379-386}). PQ is opt-in. &
\yes. Snapshots copy index files
(\f{milvus/internal/datacoord/copy_segment_task.go:803-806,990-1014});
\texttt{RestoreExternalSnapshot} (\f{milvus/internal/proxy/snapshot_impl.go:213}). &
\no\ content check. Build-ID \emph{ownership} only
(\f{milvus/internal/datacoord/copy_segment_task.go:1316-1320}). &
Not in repo (knowhere). &
\no\ public API (no centroid or codebook symbol in \f{milvus/internal/proxy}). Only by
crafting index files. \\
\midrule
Qdrant &
\yes. PQ (\f{qdrant/lib/quantization/src/encoded_vectors_pq.rs:342-400}), k-means
(\f{qdrant/lib/quantization/src/kmeans.rs:9-27}), SQ, BQ, TQ. &
\textbf{opt-in}: \texttt{quantization: null} (\f{qdrant/config/config.yaml:250-252}). HNSW is
uncompressed. Quantization is built only for indexed segments (threshold
\f{qdrant/lib/shard/src/optimizers/config.rs:17}). &
\yes. The snapshot tar includes the quantized files
(\f{qdrant/lib/segment/src/segment/snapshot.rs:304-312}); recover from a URL
(\f{qdrant/src/actix/api/snapshot_api.rs:251}). &
\no. The files are loaded as-is
(\f{qdrant/lib/segment/src/vector_storage/quantized/quantized_vectors/load.rs:32-35}); an
optional user-supplied SHA-256 covers the whole file. \textbf{PQ/SQ do not rescore by default}
(\f{qdrant/lib/segment/src/vector_storage/quantized/quantized_vectors/accessors.rs:16-22}). &
\textbf{unseeded}: \texttt{rand::rng()} for sampling
(\f{qdrant/lib/quantization/src/encoded_vectors_pq.rs:365-366}). Nothing recorded. &
\no\ API. \texttt{quantized.meta.json} is plain JSON
(\f{qdrant/lib/quantization/src/encoded_vectors_pq.rs:47-51}), so it can be edited in the
snapshot. \\
\midrule
Weaviate &
\yes. PQ (\f{weaviate/adapters/repos/db/vector/compressionhelpers/product_quantization.go:401-457}),
k-means encoder, SQ, RQ, BQ. &
\textbf{opt-in}. Default index is HNSW (\f{weaviate/entities/vectorindex/config.go:25}); PQ,
SQ, RQ, and BQ are all \texttt{false} (\f{weaviate/entities/vectorindex/hnsw/pq_config.go:28});
\texttt{DEFAULT\_QUANTIZATION} defaults to empty (\f{weaviate/usecases/config/environment.go:633-638}).
Once enabled, PQ trains automatically at 100k (\f{weaviate/adapters/repos/db/vector_index_queue.go:285-302}). &
\yes. The backup includes HNSW commit logs
(\f{weaviate/adapters/repos/db/vector/hnsw/backup.go:77-79}), which carry the PQ encoders
(\f{weaviate/adapters/repos/db/vector/hnsw/commitlog/logger.go:84-103}). &
\no\ check, \textbf{but rescores by default} with full vectors
(\f{weaviate/adapters/repos/db/vector/hnsw/search.go:188-195,822-852}). &
\textbf{random}: \texttt{Seed: rand.Uint64()} (\f{weaviate/adapters/repos/db/vector/kmeans/kmeans.go:163}).
Not persisted (\f{weaviate/entities/vectorindex/compression/pq_data.go:41-50}). &
\partl. Internal Go constructors take centers
(\f{.../compressionhelpers/kmeans_encoder.go:41-45}). There is no REST path, so the attacker
has to craft commit logs. \\
\bottomrule
\end{xltabular}}
\end{landscape}

%==========================================================================
\section{Per-repository evidence}
%==========================================================================

\subsection{FAISS}

\paragraph{Q1 (PQ construction).} \f{faiss/faiss/impl/ProductQuantizer.h:24-29} documents
the construction: ``\emph{PQ is trained using k-means, minimizing the L2 distance to
centroids.}'' The centroid table is a public \texttt{std::vector<float> centroids} with layout
$(M, k_{sub}, d_{sub})$ (\f{faiss/faiss/impl/ProductQuantizer.h:54-56}). Training slices each
sub-vector and runs an independent \texttt{Clustering} for each $m$:
\begin{lstlisting}
// faiss/faiss/impl/ProductQuantizer.cpp:143-189
for (size_t m = 0; m < M; m++) {
    ... memcpy(xslice.get() + j * dsub, x + j * d + m * dsub, ...);
    Clustering clus(dsub, ksub, cp);
    ...
    clus.train(n, xslice.get(), assign_index ? *assign_index : index);
    set_params(clus.centroids.data(), m);
}
\end{lstlisting}
Encoding assigns each sub-vector to its nearest centroid in \texttt{compute\_code}
(\f{faiss/faiss/impl/ProductQuantizer.cpp:282}). \texttt{IndexIVFPQ} (\f{faiss/faiss/IndexIVFPQ.h:34})
trains a coarse quantizer and then trains the PQ on residuals: \texttt{train\_encoder} calls
\texttt{pq.train(n, x)} (\f{faiss/faiss/IndexIVFPQ.cpp:76-80}). \texttt{by\_residual = true}
is the default (\f{faiss/faiss/IndexIVFPQ.cpp:64}). The coarse k-means uses
\texttt{cp.niter = 10} (\f{faiss/faiss/IndexIVF.cpp:45}). OPQ is \texttt{OPQMatrix}
(\f{faiss/faiss/VectorTransform.h:258-262}), trained in
\f{faiss/faiss/VectorTransform.cpp:1246}. SQ types are parsed from factory strings
(\f{faiss/faiss/index_factory.cpp:162-181}).

\paragraph{Q2 (default).} FAISS has no default index. \texttt{index\_factory(d, description, metric)}
requires a description (\f{faiss/faiss/index_factory.h:17-20}), and \texttt{"Flat"} maps to an
exhaustive index:
\begin{lstlisting}
// faiss/faiss/index_factory.cpp:808-811
// IndexFlat
if (description == "Flat") { return new IndexFlat(d, metric); }
\end{lstlisting}
\texttt{IndexFlat} ``\emph{stores the full vectors and performs exhaustive search}''
(\f{faiss/faiss/IndexFlat.h:20-21}). PQ appears only when the user writes it into the
description: \texttt{"PQ\{M\}x\{nbits\}"} (\f{faiss/faiss/index_factory.cpp:864-869}),
\texttt{"IVF\{n\},PQ\{M\}"} (\f{faiss/faiss/index_factory.cpp:420-425}), or
\texttt{"OPQ\{M\}"} (\f{faiss/faiss/index_factory.cpp:255-258}). \textbf{PQ is never the
default.}

\paragraph{Q3 (artifact).} The PQ is serialized together with its centroids:
\begin{lstlisting}
// faiss/faiss/impl/index_write.cpp:185-190
void write_ProductQuantizer(const ProductQuantizer* pq, IOWriter* f) {
    WRITE1(pq->d); WRITE1(pq->M); WRITE1(pq->nbits);
    WRITEVECTOR(pq->centroids);
}
\end{lstlisting}
For \texttt{IndexIVFPQ}, the writer emits the \texttt{IwPQ} tag, the IVF header (which contains
the coarse quantizer), the PQ, and the inverted lists of codes
(\f{faiss/faiss/impl/index_write.cpp:801-810}). The raw float vectors are not written; only
the codes are, so the full-precision embeddings are \emph{not} in the artifact. Python exposes
\texttt{serialize\_index} and \texttt{deserialize\_index} for byte buffers
(\f{faiss/faiss/python/__init__.py:564,571}).

\paragraph{Q4 (verification).} \texttt{read\_ProductQuantizer} checks structure only
(\texttt{M>0}, \texttt{nbits} in $[1,24]$, allocation limits, and
\texttt{centroids.size()==d*ksub}) (\f{faiss/faiss/impl/index_read.cpp:735-768}). SQ reads
check the length of \texttt{trained} (\f{faiss/faiss/impl/index_read.cpp:1060-1165}). A
search of \f{faiss/faiss/impl/index_read.cpp}, \f{index_write.cpp}, \f{io.cpp}, and
\f{index_io.h} for checksum, CRC, hash, verify, signature, and integrity found no integrity
mechanism. The only hardening is resource limits against malformed files
(\f{faiss/faiss/index_io.h:120-137}). Exact re-ranking is available only through an explicit
\texttt{IndexRefine} wrapper (\f{faiss/faiss/IndexRefine.h:24-39}).

\paragraph{Q5 (seed).} \texttt{ClusteringParameters::seed = 1234}. A negative seed selects a
clock-based seed (\f{faiss/faiss/Clustering.h:47-50}; \f{faiss/faiss/impl/ClusteringHelpers.cpp:28-34}).
Random initialization uses \texttt{rand\_perm(..., actual\_seed + 1 + redo*15486557L)}
(\f{faiss/faiss/Clustering.cpp:192,211}). OPQ's initial rotation uses
\texttt{float\_randn(rotation, d*d, 1234)} (\f{faiss/faiss/VectorTransform.cpp:1262,1313}).
\textbf{Consequence:} with default parameters, FAISS PQ training is reproducible given the
same training set, order, and parameters. However, \texttt{cp} (seed, \texttt{niter},
\texttt{nredo}) is \emph{not} serialized: \texttt{write\_ProductQuantizer} writes only
\texttt{d, M, nbits, centroids}, and \texttt{index\_write.cpp} contains no
\texttt{ClusteringParameters} or \texttt{cp.} reference. A verifier therefore has to guess the
parameters and needs the exact training sample. We did not check whether results are
bit-identical across thread counts or SIMD levels.

\paragraph{Q6 (direct override).} Centroids can be written directly:
\begin{lstlisting}
// faiss/faiss/impl/ProductQuantizer.cpp:76-80
void ProductQuantizer::set_params(const float* centroids_, int m) {
    memcpy(get_centroids(m, 0), centroids_, ksub * dsub * sizeof(centroids_[0]));
}
\end{lstlisting}
Other public routes: \texttt{Train\_hot\_start} uses existing centroids as the k-means
initialization (\f{faiss/faiss/impl/ProductQuantizer.h:42}; \f{.cpp:175-178});
\texttt{frozen\_centroids} keeps the supplied centroids fixed during training
(\f{faiss/faiss/Clustering.h:38-40}); Python \texttt{copy\_array\_to\_vector} writes arrays
into \texttt{pq.centroids} (\f{faiss/faiss/python/array_conversions.py:135}), as in
\f{faiss/tests/test_index_binary.py:47}. Codes can be written directly as well:
\texttt{Index::add\_sa\_codes} (\f{faiss/faiss/Index.h:430}; IVF implementation
\f{faiss/faiss/IndexIVF.cpp:213}), and \texttt{InvertedLists::update\_entry/update\_entries}
(\f{faiss/faiss/invlists/InvertedLists.h:138-149}). \textbf{The attack does not need to go
through k-means.}

\subsection{LangChain}

\paragraph{Q1.} LangChain implements no PQ or SQ. Searching \f{langchain/libs} for
\texttt{quantiz}, \texttt{ProductQuant}, and \texttt{ScalarQuant} outside tests finds only
Qdrant pass-through parameters and one unrelated embedding comment. The FAISS wrapper is a
lazy re-export from \texttt{langchain\_community}:
\begin{lstlisting}
# langchain/libs/langchain/langchain_classic/vectorstores/faiss.py:11
DEPRECATED_LOOKUP = {"FAISS": "langchain_community.vectorstores"}
\end{lstlisting}
The same pattern appears in \f{langchain/libs/langchain/langchain_classic/__init__.py:348-353}.
\textbf{\texttt{langchain\_community} is not in this repository}, so we could not audit
LangChain's FAISS default index or its \texttt{load\_local}/\texttt{save\_local} behavior from
the cloned code.

\paragraph{Q2.} \texttt{InMemoryVectorStore} (\f{langchain/libs/core/langchain_core/vectorstores/in_memory.py:34})
computes exact cosine over every stored vector
(\f{langchain/libs/core/langchain_core/vectorstores/in_memory.py:314}). The
\texttt{langchain\_qdrant.QdrantVectorStore.from\_texts} path defaults to
\texttt{distance=COSINE} (\f{langchain/libs/partners/qdrant/langchain_qdrant/qdrant.py:357})
and creates the collection with only \texttt{VectorParams(size, distance)}
(\f{langchain/libs/partners/qdrant/langchain_qdrant/qdrant.py:944-982}). The string
\texttt{quantiz} does not occur anywhere in \f{langchain_qdrant/qdrant.py}. Only the
deprecated \texttt{Qdrant} class (\f{langchain/libs/partners/qdrant/langchain_qdrant/vectorstores.py:59-60})
accepts \texttt{quantization\_config}, and its docstring says ``\emph{if None - quantization
will be disabled}'' (\f{.../vectorstores.py:1309-1310}). The Chroma partner contains no
quantization code (search of \f{langchain/libs/partners/chroma/langchain_chroma/} for
\texttt{quantiz} and \texttt{pq}: none).

\paragraph{Q3.} \texttt{InMemoryVectorStore.dump/load} writes and reads full-precision JSON
(\f{langchain/libs/core/langchain_core/vectorstores/in_memory.py:517-546}), so nothing is
compressed. \texttt{QdrantVectorStore.from\_existing\_collection} attaches to a collection that
someone else built (\f{langchain/libs/partners/qdrant/langchain_qdrant/qdrant.py:434-470}).
That collection could be quantized, but LangChain does not create the quantization.

\paragraph{Q4.} When the collection already exists, \texttt{\_validate\_collection\_config}
checks only vector size (\f{.../qdrant.py:1230}) and distance (\f{.../qdrant.py:1240}). It
reads no quantization configuration and compares nothing to source embeddings.

\paragraph{Q5/Q6.} Not applicable in this repository.

\subsection{LlamaIndex}

\paragraph{Q1.} LlamaIndex has no own PQ. \texttt{FaissVectorStore} takes a caller-constructed
\texttt{faiss\_index} (\f{.../llama-index-vector-stores-faiss/llama_index/vector_stores/faiss/base.py:65-82};
below, \f{.../faiss/} abbreviates that integration directory).

\paragraph{Q2.} With no vector store supplied, \texttt{StorageContext.from\_defaults} creates
\texttt{SimpleVectorStore()}
(\f{llama_index/llama-index-core/llama_index/core/storage/storage_context.py:105-107}). It
holds \texttt{Dict[str, List[float]]} embeddings
(\f{llama_index/llama-index-core/llama_index/core/vector_stores/simple.py:48-62}) and runs
exact \texttt{get\_top\_k\_embeddings} (\f{.../simple.py:303}). The FAISS integration has no
default index type: the user must pass one. The docstring and every doc example use
\texttt{faiss.IndexFlatL2(d)} (\f{.../faiss/base.py:54};
\f{llama_index/docs/examples/vector_stores/FaissIndexDemo.ipynb:84};
\f{llama_index/docs/src/content/docs/framework/community/integrations/vector_stores.md:487,495-496}).
A search of \f{llama_index/docs} for \texttt{IndexIVFPQ} and \texttt{faiss.IndexPQ} found no
use. The Milvus integration overrides Milvus's server default and asks for \texttt{FLAT}:
\texttt{index\_type = base\_params.pop("index\_type", "FLAT")}
(\f{.../llama-index-vector-stores-milvus/llama_index/vector_stores/milvus/base.py:1406}).
The Qdrant integration creates \texttt{VectorParams(size, COSINE)} and forwards
\texttt{quantization\_config}, which defaults to \texttt{None}
(\f{.../llama-index-vector-stores-qdrant/llama_index/vector_stores/qdrant/base.py:190,810-815,835}).
The Weaviate integration's default schema sets no vector-index configuration
(\f{.../llama-index-vector-stores-weaviate/llama_index/vector_stores/weaviate/utils.py:100-108}).

\paragraph{Q3 (the strongest artifact boundary we found).} \texttt{persist} calls
\texttt{faiss.write\_index(self.\_faiss\_index, persist\_path)} (\f{.../faiss/base.py:171}),
and \texttt{from\_persist\_dir}/\texttt{from\_persist\_path} read it back
(\f{.../faiss/base.py:85-117}). The core index also \emph{removes} embeddings before writing
nodes to the docstore when the vector store does not store text (FAISS sets
\texttt{stores\_text = False}, \f{.../faiss/base.py:61}):
\begin{lstlisting}
# llama_index/llama-index-core/llama_index/core/indices/vector_store/base.py:196-199
# NOTE: remove embedding from node to avoid duplication
node_without_embedding = node.model_copy()
node_without_embedding.embedding = None
\end{lstlisting}
The same code appears in the sync path at lines 237--240. In a persisted LlamaIndex+FAISS
directory, the FAISS file is therefore the \emph{only} copy of the vectors. With IVFPQ, no
full-precision embedding remains anywhere in the artifact.

\paragraph{Q4.} There is no verification:
\begin{lstlisting}
# .../faiss/base.py:113-117
logger.info(f"Loading {__name__} from {persist_path}.")
faiss_index = faiss.read_index(persist_path)
return cls(faiss_index=faiss_index)
\end{lstlisting}
The index type is not checked. \texttt{FaissVectorMapStore} checks only that the loaded object
is an \texttt{IndexIDMap}/\texttt{IndexIDMap2} (\f{.../faiss/map_store.py:89-94}); it does
not check what is inside.

\paragraph{Q5/Q6.} These are inherited from FAISS. Any FAISS object, including one with
hand-set centroids, is accepted.

\subsection{Milvus}

\paragraph{Q1.} The Go SDK exposes quantized index types: \texttt{IVF\_PQ},
\texttt{IVF\_SQ8}, \texttt{HNSW\_SQ}, \texttt{HNSW\_PQ}, and \texttt{SCANN}
(\f{milvus/client/index/common.go:46-57}). The implementation is in \texttt{knowhere}, which is
fetched at build time from \texttt{https://github.com/zilliztech/knowhere.git} at version
\texttt{faff72c} (\f{milvus/internal/core/thirdparty/knowhere/CMakeLists.txt:17-18}). Index
builds go through cgo (\f{milvus/internal/util/indexcgowrapper/index.go:109-119}).
\textbf{The PQ/SQ training code is not in this repository}, so we cannot cite it.

\paragraph{Q2 (the one default-on case).} The CPU-build default for dense float vectors is:
\begin{lstlisting}
// milvus/pkg/util/paramtable/autoindex_param_nocuda.go:26
DefaultValue: `{"M": 18 ,"efConstruction": 240,"index_type": "HNSW_SQ",
  "sq_type": "SQ4U", "refine": true, "refine_type": "FP16", "metric_type": "COSINE"}`,
\end{lstlisting}
The CUDA build uses \texttt{GPU\_CAGRA} (\f{milvus/pkg/util/paramtable/autoindex_param_cuda.go:26}).
The \texttt{largeTopK} variant is \texttt{IVF\_SQ8} (\f{milvus/pkg/util/paramtable/autoindex_param.go:155}),
but it applies only when \texttt{autoIndex.enable} is true
(\f{milvus/internal/util/indexparamcheck/index_params_validation.go:181-187}), and that
defaults to \texttt{false} (\f{milvus/pkg/util/paramtable/autoindex_param.go:82-84}), a flag
the code comments as ``\emph{`enable` only for cloud instance}''
(\f{milvus/internal/proxy/task_index.go:249}). On open-source servers, the
``behavior change after 2.2.9'' branch applies the AUTOINDEX configuration whenever
\texttt{index\_type} is absent or equals \texttt{AUTOINDEX}
(\f{milvus/internal/proxy/task_index.go:311,379-386}). The Go SDK's quick setup,
\texttt{SimpleCreateCollectionOptions} (\f{milvus/client/milvusclient/collection_options.go:189-205}),
creates an index with an empty parameter map:
\begin{lstlisting}
// milvus/client/milvusclient/collection_options.go:176-179
if opt.isFast && opt.indexOptions == nil {
    return []CreateIndexOption{ NewCreateIndexOption(opt.name, opt.vectorFieldName,
        index.NewGenericIndex("", map[string]string{})) } }
\end{lstlisting}
The client issues it (\f{milvus/client/milvusclient/collection.go:52-54}), so quick setup
lands on \texttt{HNSW\_SQ}/\texttt{SQ4U}. This is \emph{scalar} quantization (4-bit uniform
ranges), not a PQ codebook. PQ variants are opt-in by name. We did not verify pymilvus's
quick-setup path, which is not in the repo.

\paragraph{Q3.} Milvus keeps raw vectors in binlogs and index files in separate objects.
\texttt{CreateSnapshot} and \texttt{RestoreSnapshot} (\f{milvus/internal/proxy/snapshot_impl.go:37,443})
and \texttt{RestoreExternalSnapshot} (\f{milvus/internal/proxy/snapshot_impl.go:213}; datacoord
\f{milvus/internal/datacoord/ddl_callbacks_snapshot.go:132-143}) restore through copy-segment
jobs. The code says ``\emph{Index files (vector/scalar/text/JSON) need to be copied with segment
data}'' (\f{milvus/internal/datacoord/copy_segment_task.go:805}) and ``\emph{WHICH INDEX FILES
TO COPY: the snapshot's own index metadata wins}'' (\f{.../copy_segment_task.go:990}). Built
indexes, including quantizer parameters and refine data, therefore travel as files and are
not rebuilt.

\paragraph{Q4.} The only post-copy check is \texttt{verifyCopiedManifestIndexOwnership}. It
``\emph{binds each published build to its actual target index ID}''
(\f{milvus/internal/datacoord/copy_segment_task.go:1316-1320}): an identity and path check,
not a content check. We found no recomputation of index contents from binlogs on restore.

\paragraph{Q5.} Not determinable from this repository (knowhere).

\paragraph{Q6.} No public API. A search of \f{milvus/internal/proxy/*.go} (non-test) for
\texttt{centroid} and \texttt{codebook} returns nothing. An attacker has to craft index files,
for example inside an external snapshot.

\subsection{Qdrant}

\paragraph{Q1.} The PQ encoder splits vectors into chunks, samples up to
\texttt{KMEANS\_SAMPLE\_SIZE = 10\_000} vectors, and runs k-means with
\texttt{CENTROIDS\_COUNT = 256} per chunk
(\f{qdrant/lib/quantization/src/encoded_vectors_pq.rs:28-31,342-400}). k-means initializes
from the first $k$ sampled points (\f{qdrant/lib/quantization/src/kmeans.rs:26-27}). The
centroids are stored in \texttt{Metadata \{ centroids: Vec<Vec<f32>>, vector\_division,
vector\_parameters \}} (\f{.../encoded_vectors_pq.rs:47-51}). Compression ratios run from
\texttt{X4} to \texttt{X64} (\f{qdrant/lib/segment/src/types.rs:922-928}). SQ is int8 with an
optional quantile (\f{qdrant/lib/segment/src/types.rs:939-957};
\f{qdrant/lib/quantization/src/encoded_vectors_u8.rs:193-205}).

\paragraph{Q2.} Quantization is off by default. The collection parameter is
\texttt{\#[serde(default)] quantization\_config: Option<QuantizationConfig>}
(\f{qdrant/lib/collection/src/config.rs:342-344}), and the server default is:
\begin{lstlisting}
# qdrant/config/config.yaml:250-252
# Default quantization configuration.
quantization: null
\end{lstlisting}
When quantization is configured, it is built only for indexed/optimized segments. Appendable
segments get quantization only under a feature flag
(\f{qdrant/lib/segment/src/types.rs:1133-1137}; \f{qdrant/lib/shard/src/optimizers/config.rs:44}).
Segments are indexed once they pass \texttt{DEFAULT\_INDEXING\_THRESHOLD\_KB = 10\_000}
(\f{qdrant/lib/shard/src/optimizers/config.rs:17}; \f{qdrant/config/config.yaml:169}). This
is a size trigger.

\paragraph{Q3.} Snapshots add every quantized-vector file to the tar
(\f{qdrant/lib/segment/src/segment/snapshot.rs:304-312}). The file names are
\texttt{quantized.config.json}, \texttt{quantized.data}, and \texttt{quantized.meta.json}
(\f{qdrant/lib/segment/src/vector_storage/quantized/quantized_vectors/config.rs:10-14}).
Snapshots can be uploaded or recovered from a location (\f{qdrant/src/actix/api/snapshot_api.rs:193,251}),
including remote URLs with an optional API key
(\f{qdrant/lib/collection/src/operations/snapshot_ops.rs:82}). The original vectors are also
in the snapshot, because vector storage files are appended just before
(\f{.../snapshot.rs:295-302}), so verification would be \emph{possible}.

\paragraph{Q4.} When a quantization configuration file exists, it is loaded without
recomputation:
\begin{lstlisting}
// qdrant/lib/segment/src/vector_storage/quantized/quantized_vectors/load.rs:32-35
let config_path = Self::get_config_path(path);
if config_path.exists() {
    let config: QuantizedVectorsConfig = read_json(&config_path)?;
    return Ok(Some(Self::load_impl(config, vector_storage, path)?));
}
\end{lstlisting}
\texttt{EncodedVectorsPQ::load} validates only the storage vector size against the metadata
(\f{qdrant/lib/quantization/src/encoded_vectors_pq.rs:144-162}). The snapshot checksum is
``\emph{Optional SHA256 checksum to verify snapshot integrity}''
(\f{qdrant/lib/collection/src/operations/snapshot_ops.rs:77-80};
\f{qdrant/lib/storage/src/content_manager/snapshots/recover.rs:136-146}). It is optional and
comes from the same party as the file. \textbf{Rescoring:} PQ and SQ default to no rescoring:
\begin{lstlisting}
// .../quantized/quantized_vectors/accessors.rs:16-22
QuantizedVectorStorage::ScalarRam(_) => false, ... PQRam(_) => false, ...
QuantizedVectorStorage::BinaryRam(_) => true,
\end{lstlisting}
The effective choice is \texttt{params...rescore.unwrap\_or(default\_rescoring)}
(\f{qdrant/lib/segment/src/index/vector_index_search_common.rs:64-72}), and oversampling
defaults to 1.0 (\f{qdrant/lib/segment/src/types.rs:588-597}). With PQ enabled, the compressed
scores are the final scores unless the query asks for rescoring. A consumer who changes the
collection's quantization configuration forces a rebuild through the config-mismatch
optimizer (\f{qdrant/lib/shard/src/optimizers/config_mismatch_optimizer.rs:110-141}).

\paragraph{Q5.} PQ sampling uses a thread-local OS-seeded RNG:
\texttt{rand::seq::index::sample(\&mut rand::rng(), count, sample\_size)}
(\f{qdrant/lib/quantization/src/encoded_vectors_pq.rs:365-366}). Empty clusters are reseeded
with \texttt{rand::rng()} (\f{qdrant/lib/quantization/src/kmeans.rs:114-115}), and the SQ
quantile sample also uses \texttt{rand::rng()} (\f{qdrant/lib/quantization/src/quantile.rs:290}).
No seed is stored in \texttt{Metadata}. \textbf{Qdrant codebooks cannot be reproduced by
recomputation}, so a malicious codebook cannot be told apart from an unlucky honest one that
way.

\paragraph{Q6.} Searching \f{qdrant/lib/api} and \f{qdrant/src} for \texttt{centroid} and
\texttt{codebook} returns nothing, so there is no public API. Centroids are plain JSON in
\texttt{quantized.meta.json}, and editing a snapshot is trivial.

\subsection{Weaviate}

\paragraph{Q1.} \texttt{ProductQuantizer.Fit} truncates the training data to
\texttt{trainingLimit} and fits one \texttt{KMeansEncoder} per segment
(\f{weaviate/adapters/repos/db/vector/compressionhelpers/product_quantization.go:401-457}).
\texttt{Encode} maps each segment to its nearest centroid (\f{.../product_quantization.go:460-466}).
The encoder uses random initialization (\f{.../compressionhelpers/kmeans_encoder.go:48-65}).
SQ, RQ, and BQ live in the same package.

\paragraph{Q2.} The default vector index is \texttt{HNSW} (\f{weaviate/entities/vectorindex/config.go:25}).
Every compressor is off by default: \texttt{DefaultPQEnabled = false}
(\f{weaviate/entities/vectorindex/hnsw/pq_config.go:28}), \texttt{DefaultSQEnabled = false}
(\f{.../sq_config.go:17}), \texttt{DefaultRQEnabled = false} (\f{.../rq_config.go:22}), and
\texttt{DefaultBQEnabled = false} (\f{.../bq_config.go:17}). A server-wide default comes from
the environment and is empty unless set:
\begin{lstlisting}
// weaviate/usecases/config/environment.go:633-638
defaultQuantization := ""
if v := os.Getenv("DEFAULT_QUANTIZATION"); v != "" { ... }
\end{lstlisting}
\texttt{ParseDefaultQuantization} maps \texttt{"pq"}, \texttt{"sq"}, \texttt{"rq-1/4/8"}, and
\texttt{"bq"} to enabled flags; \texttt{""} and \texttt{"none"} leave compression off
(\f{weaviate/entities/vectorindex/hnsw/config.go:345-381}). In-repo documentation shows that
setting \texttt{ALLOWED\_COMPRESSION\_TYPES} to a single entry seeds
\texttt{DEFAULT\_QUANTIZATION} (example \texttt{rq-8}) (\f{weaviate/docs/usage_limits.md:111,121,130-139}).
That is an operator choice, not a default. We did not audit Weaviate Cloud defaults or the
external documentation site. \textbf{Scale trigger:} once PQ is enabled, the index queue
calls \texttt{Upgrade} as soon as \texttt{AlreadyIndexed() > shouldUpgradeAt}
(\f{weaviate/adapters/repos/db/vector_index_queue.go:285-302}). For PQ that threshold is
\texttt{pqConfig.TrainingLimit} (\f{weaviate/adapters/repos/db/vector/hnsw/index.go:976}),
with a default of 100{,}000 (\f{.../pq_config.go:34}).

\paragraph{Q3.} Backups list the HNSW commit-log directory
(\f{weaviate/adapters/repos/db/vector/hnsw/backup.go:77-79}). The commit log records the PQ
codebooks: \texttt{AddPQCompression} writes the header and then
\texttt{encoder.ExposeDataForRestore()} for every segment encoder
(\f{weaviate/adapters/repos/db/vector/hnsw/commitlog/logger.go:84-103}). On restart, the
encoders are rebuilt from these bytes through \texttt{NewProductQuantizerWithEncoders}
(\f{.../product_quantization.go:243-253}) instead of being retrained. Full-precision vectors
remain in the object store.

\paragraph{Q4.} In \f{weaviate/usecases/backup/*.go}, the only matches for checksum, verify,
and integrity are coverage and coordination checks
(\f{weaviate/usecases/backup/coordinator_dedupe.go:449-451}). There is no content
verification. \textbf{However, Weaviate rescores by default}:
\begin{lstlisting}
// weaviate/adapters/repos/db/vector/hnsw/search.go:188-195
func (h *hnsw) shouldRescore() bool {
  if h.compressed.Load() { if (h.sqConfig.Enabled && h.sqConfig.RescoreLimit == 0) || ... { return false } }
  return h.compressed.Load() && !h.doNotRescore }
\end{lstlisting}
The rescore distance fetches the full-precision vector through
\texttt{TempVectorForIDWithViewThunk} (\f{.../search.go:822-852}). \texttt{doNotRescore}
defaults to the Go zero value (\f{weaviate/adapters/repos/db/vector/hnsw/index.go:180}). A
malicious PQ codebook in Weaviate can therefore only change \emph{which candidates reach the
rescoring stage}, not their final order.

\paragraph{Q5.} \texttt{kmeans.New} sets \texttt{Seed: rand.Uint64()}
(\f{weaviate/adapters/repos/db/vector/kmeans/kmeans.go:163}), which feeds a PCG
(\f{.../kmeans.go:50}). The training sample comes from a sparse Fisher--Yates sampler that
uses global \texttt{rand.Intn} (\f{weaviate/adapters/repos/db/vector/common/iterator.go:46};
used in \f{weaviate/adapters/repos/db/vector/hnsw/compress.go:48-75}). \texttt{PQData} has no
seed field (\f{weaviate/entities/vectorindex/compression/pq_data.go:41-50}). Training is not
reproducible.

\paragraph{Q6.} \texttt{NewKMeansEncoderWithCenters(k, d, segment, centers)}
(\f{.../compressionhelpers/kmeans_encoder.go:41-45}) and \texttt{NewProductQuantizerWithEncoders}
are internal Go APIs. \f{weaviate/adapters/handlers/rest/operations} has no REST operation that
references centroids or codebooks. An attacker would have to craft commit logs inside a backup.

%==========================================================================
\section{Threats to the attack premise}
%==========================================================================
\begin{enumerate}
\item \textbf{PQ is opt-in in every repository.} No code path we traced turns on product
  quantization without the user naming it: a FAISS factory string, a Qdrant
  \texttt{quantization\_config}, Weaviate \texttt{pq.enabled} or
  \texttt{DEFAULT\_QUANTIZATION=pq}, or Milvus \texttt{index\_type=IVF\_PQ/HNSW\_PQ}.
\item \textbf{The one default-on quantizer (Milvus) is SQ with refine.} It uses 4-bit uniform
  SQ plus an FP16 refine stage (\f{milvus/pkg/util/paramtable/autoindex_param_nocuda.go:26}).
  SQ has far fewer degrees of freedom than a PQ codebook: in effect, per-dimension ranges. This
  matches our prior finding that calibration ranges, not weights, are the malicious degree of
  freedom. Refine should re-rank candidates at FP16, which limits the attack to suppression.
  We could not verify the refine semantics because they live in knowhere, outside the repo.
\item \textbf{Default rescoring in Weaviate} with full-precision vectors
  (\f{weaviate/adapters/repos/db/vector/hnsw/search.go:188-195,822-852}) blocks promotion. A
  document that is not truly close cannot enter the final top-$k$ once rescored, because only
  candidate recall can be manipulated.
\item \textbf{Scale and index gating.} Qdrant quantizes only indexed segments
  ($\ge$10{,}000\,KB, \f{qdrant/lib/shard/src/optimizers/config.rs:17}). Weaviate PQ trains
  only after 100{,}000 indexed vectors (\f{weaviate/entities/vectorindex/hnsw/pq_config.go:34}).
  Small RAG corpora may never reach the quantized path.
\item \textbf{Framework defaults are exact.} LangChain \texttt{InMemoryVectorStore} and
  LlamaIndex \texttt{SimpleVectorStore} are exact and full precision. The LlamaIndex Milvus
  integration even forces \texttt{FLAT} over Milvus's own SQ default
  (\f{.../milvus/base.py:1406}).
\item \textbf{Rebuild on configuration change (Qdrant).} A consumer who restores a snapshot and
  then sets their own quantization configuration triggers a rebuild that discards the shipped
  codebook (\f{qdrant/lib/shard/src/optimizers/config_mismatch_optimizer.rs:110-141}).
\item \textbf{Verification is possible in the vector DBs.} Qdrant snapshots, Weaviate's object
  store, and Milvus binlogs all keep the original vectors next to the compressed index. A
  consumer \emph{could} re-encode them and compare. None of them does so, but the data needed
  for detection is present. This is not true of FAISS IVFPQ or LlamaIndex+FAISS.
\item \textbf{FAISS training is deterministic by default} (\texttt{seed=1234}). With the training
  set and parameters, an honest codebook can be recomputed exactly, and a malicious one would
  differ. The parameters are not serialized, however, and FAISS IVFPQ files contain no training
  vectors.
\item \textbf{The attacker who ships the artifact has more power than the attack needs.} In
  every Q3 path, the party shipping the file can edit codes or ID mappings directly
  (FAISS \texttt{add\_sa\_codes} and \texttt{update\_entries}; Qdrant JSON metadata). The
  quantization-conditioned framing adds value only against a consumer who spot-checks exact
  behavior, or who checks that the codes are consistent with the shipped codebook. No repo
  performs either check today. This should be stated explicitly in the threat model.
\item \textbf{Unaudited code outside the repositories.} LangChain's FAISS wrapper
  (\texttt{langchain\_community}), Milvus's quantizers (knowhere), pymilvus quick setup, and
  hosted-service defaults (Zilliz, Weaviate Cloud, Qdrant Cloud) are not in these repositories
  and were not audited.
\end{enumerate}

%==========================================================================
\section{Strongest instantiation}\label{sec:strongest}
%==========================================================================
\textbf{Repo:} LlamaIndex + FAISS. \textbf{Index type:} \texttt{IndexIVFPQ}, optionally with
\texttt{OPQ}, e.g.\ \texttt{faiss.index\_factory(d, "OPQ16,IVF1024,PQ16")}.
\textbf{Configuration:} the documented persist-and-reload pattern.

Why this one:
\begin{enumerate}
\item \emph{The shipped artifact is the only copy of the vectors.} IVFPQ stores codes only
  (\f{faiss/faiss/impl/index_write.cpp:801-810}), and LlamaIndex removes embeddings from the
  docstore (\f{llama_index/llama-index-core/llama_index/core/indices/vector_store/base.py:196-199}).
  The consumer has nothing to verify against unless they re-embed the corpus.
\item \emph{Loading does not depend on the index type.} \texttt{from\_persist\_path} accepts
  whatever \texttt{faiss.read\_index} returns (\f{.../faiss/base.py:113-117}). A consumer who
  followed the \texttt{IndexFlatL2} example receives a PQ index silently.
\item \emph{There is no rescoring.} IVFPQ search ranks by PQ-approximated distance unless the
  builder wrapped it in \texttt{IndexRefine}. Promotion and suppression both work.
\item \emph{Codebooks and codes can be written directly}
  (\f{faiss/faiss/impl/ProductQuantizer.cpp:76-80}; \f{faiss/faiss/Index.h:430}). The
  attacker does not need to steer k-means.
\item \emph{Training parameters are not recorded} in the file, so ``re-train and compare''
  requires guessing the parameters.
\end{enumerate}

\noindent\textbf{Consumer steps that land on a PQ index they did not build} (the pattern in
\f{llama_index/docs/examples/vector_stores/FaissIndexDemo.ipynb:177-181}):
\begin{lstlisting}[language=Python]
# 1. Obtain ./storage from a third party (teammate, data vendor, model/dataset hub, CI artifact).
#    It contains default__vector_store.json (a FAISS binary written by faiss.write_index),
#    docstore.json (nodes WITHOUT embeddings), index_store.json.
from llama_index.vector_stores.faiss import FaissVectorStore
from llama_index.core import StorageContext, load_index_from_storage
vector_store = FaissVectorStore.from_persist_dir("./storage")      # faiss.read_index, no checks
storage_context = StorageContext.from_defaults(vector_store=vector_store, persist_dir="./storage")
index = load_index_from_storage(storage_context=storage_context)
# 2. Query. Retrieval runs on the attacker's PQ codebook/codes; the docstore maps
#    positional FAISS ids -> node text. No step inspects index type, centroids, or codes.
retriever = index.as_retriever(similarity_top_k=5)
\end{lstlisting}

\noindent\textbf{Runner-up instantiations.}
(i)~\emph{Qdrant with PQ enabled}: the consumer recovers a collection snapshot from a URL
(\f{qdrant/src/actix/api/snapshot_api.rs:251}). \texttt{quantized.meta.json} is loaded as-is
(\f{.../quantized_vectors/load.rs:32-35}), and PQ queries do not rescore by default
(\f{.../accessors.rs:16-22}). This case is weaker than FAISS because the original vectors
ship alongside the codes, so verification is possible, and because a configuration change
forces a rebuild.
(ii)~\emph{Milvus default}: any collection whose vector index was created without
\texttt{index\_type} gets \texttt{HNSW\_SQ}/\texttt{SQ4U}, and
\texttt{RestoreExternalSnapshot} copies prebuilt index files with only an ownership check.
Quantization is on by default here, but it is SQ with FP16 refine, and the quantizer code is
outside the audited repository.

\vspace{0.5em}
\noindent\textbf{Implication for the threat model.} The honest framing is: \emph{``quantized
indexes are opt-in, or scale-triggered once opted in, everywhere except Milvus's SQ default.
Once a compressed index exists, it is a trusted, shippable artifact that no audited system
verifies.''} The attack targets \emph{configured} deployments at scale and pre-built index
distribution. It does not target a default-on PQ supply chain.

\end{document}
